\section{The Gaussian operator path}
\label{S:actual-operator-driver}

We now construct the integrated operator path in the interaction representation
on the same probability space as the original noise. Retain
Hypothesis~\ref{H:homog}, set $a = \alpha_0$, $s = 1/2$ and $E_r =
H^r(\R^d)\times H^{r-1}(\R^d)$, and initially use real phases $E = E_s$. Let
$H_W$ be the real completion of $C_c^\infty((0,\infty)\times\R^d)$ modulo its
covariance-null tests. The given isonormal process on it is $W$. If $\Gamma =
0$, the operator path and all covariance forms are zero, so the conclusions
below hold directly. Assume $\Gamma \ne 0$ in the construction.

Choose and fix $R,T > 0$ and a real even nonnegative $\zeta \in
C_c^\infty(\R^d)$ equal to one on $B_1$ and zero off $B_2$. The physical mask is
$\zeta_R(x) = \zeta(x/R)$. The free group $S(t)$ and forms $\ell_R$ are those of
Section~\ref{S:wave-variance}; both phase entries and the continuous
zero-frequency sine quotient are retained. Fix parameters
\begin{equation}\label{E:driver-parameters}
  \LeanEqMarker{E:driver-parameters}
  0<\gamma<\min\left(\frac12,\frac{2-\alpha}{2}\right),\qquad
  \frac12<\theta<\beta<\min\left(1-\frac a2,
  \frac{3-a-\alpha}{2}\right).
\end{equation}
The strict intervals are nonempty. The $C^\theta$ norm is the supremum norm plus
the $\theta$-H\"older seminorm.

\begin{proposition}[Compact operator path on the original noise space]\label{P:actual-operator-driver}
  There is an orthonormal basis $\{f_j\}$ of $H_W$ with real smooth compact test
  representatives, independent standard real Gaussians $g_j = W(f_j)$, and
  deterministic operators $A_j(0,t)$ on $E$ such that
  \begin{equation}\label{E:driver-path-convergence}
    X_N(t):=\sum_{j\le N}g_jA_j(0,t)\longrightarrow M_R(t)
    \quad\text{in }L^p(\Omega;C^\theta([0,T];\mathcal{B}(E)))
  \end{equation}
  for every finite $p \ge 1$. The limit is a centered Gaussian path on the
  original probability space, starts at zero, and has compact values at every
  time on one probability-one path set. For each fixed real Schwartz pair $z,w$
  and $t \in [0,T]$,
  \begin{equation}\label{E:driver-raw-row}
    \langle M_R(t)z,w\rangle_E=\ell_R((0,t);z,w)
    \quad\text{in }L^2(\Omega).
  \end{equation}
  For long intervals, use the completed covariance test, equivalently the
  additive unit-interval construction justified below. Different compact smooth
  test bases give indistinguishable paths on the same $W$. For constants $C,c >
  0$ depending on the fixed data but not the basis or $p$,
  \begin{equation}\label{E:driver-gaussian-norm}
    (\E\|M_R\|_{C^\theta}^p)^{1/p}\le C\sqrt p,\qquad
    \E\exp(c\|M_R\|_{C^\theta}^2)<\infty.
  \end{equation}
  Its complexification has the corresponding sesquilinear covariance identity.
  Fixed-interval Schatten moments hold at the admissible even indices
  in~\eqref{E:driver-fixed-schatten}. Arbitrary $p$ above refers to
  $\mathcal{B}(E)$, not arbitrary Schatten indices or a Schatten-valued path.
\end{proposition}

\subsection{A test basis and explicit covariance coefficients}

\begin{lemma}[Basis in the prescribed quotient]\label{L:driver-noise-basis}
  The nonzero space $H_W$ is separable and infinite-dimensional. It has a real
  orthonormal basis with compact smooth representatives, whose evaluations by
  $W$ are independent standard real Gaussians.
\end{lemma}

\begin{proof}
  We first obtain a countable dense collection of tests in the prescribed
  covariance norm. We then show that the quotient is infinite-dimensional and
  apply Gram--Schmidt within that quotient.

  \medskip
  \noindent\textit{Step 1.~}We prove separability of the test completion.
  For tests supported in $[0,L]\times K$, with $K$ compact, the physical
  covariance and the temporal Schur bound give the absolute estimate
  \begin{equation}\label{E:driver-test-continuity}
    |\langle F,G\rangle_{H_W}|
    \le |K|\Gamma(K-K)\frac{2L^{1-a}}{1-a}
    \|\|F(t)\|_\infty\|_{L^2(0,L)}
    \|\|G(t)\|_\infty\|_{L^2(0,L)}.
  \end{equation}
  The spatial covariance is bounded by the product of the two spatial supremum
  norms, the volume $|K|$ and the mass $\Gamma(K-K)$. The temporal Schur
  estimate then bounds the integral of this product by the two $L^2$ time
  norms. Thus, uniform approximation on a fixed box implies approximation in
  the covariance seminorm.
  Only spatial differences in $K-K$ contribute. On each rational compact box
  inside positive time, rational step approximations followed by compact smooth
  mollifiers at rational scales give a countable uniformly dense test
  collection, supported in a fixed slightly larger box. The estimate proves
  noise-norm approximation. The union over increasing boxes is dense in the test
  completion. This proves separability without assuming the covariance form is
  injective.

  \medskip
  \noindent\textit{Step 2.~}We show that infinitely many temporal tests
  survive the quotient.
  Choose $L \ge 1$ with $0 < \Gamma(B_L) < \infty$ and nonnegative compact
  smooth $q \ge 1$ on $B_{2L+1}$. Then
  \[
    Q_{\rm sp}:=\int\Gamma(\ud y)\int q(x)q(x-y)\ud x
    \ge |B_{L+1}|\Gamma(B_L)>0,
  \]
  and the compact-support estimate makes it finite. For every real compact
  smooth positive-time $\tau$, the tensor identity is
  \[
    \|\tau\otimes q\|_{H_W}^2=Q_{\rm sp}\|\tau\|_a^2.
  \]
  For $y \in B_L$ and $x \in B_{L+1}$, both $x$ and $x-y$ lie in $B_{2L+1}$.
  Each factor in the integrand is therefore at least one on this region.
  Here the temporal norm is computed on any finite interval containing the
  support of $\tau$. It is strictly positive for signed nonzero $\tau$. Indeed,
  zero-extend this deterministic test and use the positive Laplace factorization
  from Lemma~\ref{L:masked-time-completion} to write
  \begin{align*}
    \|\tau\|_a^2
    &=\frac1{\Gamma_E(a)}\int_0^\infty\lambda^{a-1}
    E_\lambda(\tau)\ud\lambda,\\
    E_\lambda(\tau)
    &=\int_\R 2\lambda\left|\int_r^\infty
    e^{-\lambda(t-r)}\tau(t)\ud t\right|^2\ud r.
  \end{align*}
  Signed exchanges are absolutely integrable on the test support, by the
  temporal Schur estimate. If the norm vanished, some $\lambda > 0$ would have
  $E_\lambda = 0$. Its inner integral $v(r)$ is continuous and satisfies $v' =
  \lambda v-\tau$. Vanishing almost everywhere makes it identically zero and
  forces the derivative $v'$ to vanish as well. The identity $v' = \lambda
  v-\tau$ then
  forces $\tau = 0$. Thus, infinitely many linearly independent compact time
  bumps inject into $H_W$ after multiplication by $q$.

  \medskip
  \noindent\textit{Step 3.~}We construct the basis and its Gaussian
  coordinates.
  Apply Gram--Schmidt to a countable dense real test collection, discarding zero
  residuals in the given null quotient. Infinitude gives an infinite basis; each
  representative is a finite linear combination of compact smooth tests. The
  isonormal process depends only on the quotient class, so a zero residual
  really is a zero noise variable. The surviving normalized vectors have
  covariance equal to their Hilbert inner product. Thus, the
  isonormal identity gives $\E g_jg_k = \delta_{jk}$. Joint Gaussian
  characteristic functions then factor, giving independence and unit variances.
  No physical null direction is removed beyond the prescribed quotient, and no
  $H_W$-valued white Gaussian vector is used. This proves
  Lemma~\ref{L:driver-noise-basis}.
\end{proof}

The operator-norm topology is essential here: in dimensions two and three,
No-go~\ref{P:nogo-21} rules out a compatible square-integrable
Hilbert--Schmidt realization on a same-index Sobolev phase pair,
for suitable coefficient masks and a sufficiently large observation cutoff.


\begin{lemma}[Covariance coefficients]\label{L:driver-covariance-coefficients}
  For each basis representative put
  \begin{equation}\label{E:driver-covariance-coefficient}
    H_{j,R}(t,x)=\zeta_R(x)\int_{u>0}|t-u|^{-a}
    \int\Gamma(\ud y)f_j(u,x-y)\ud u.
  \end{equation}
  These functions are real smooth on finite time-space compacts, supported
  spatially in $B_{2R}$, with bounded derivatives there. For smooth compact real
  positive-time $F$, and also for the admitted wave tests
  of~\eqref{E:variance-actual-form},
  \begin{equation}\label{E:driver-literal-covariance}
    \E[W_{\mathbb{C}}(\zeta_RF)g_j]
    =\int H_{j,R}(t,x)F(t,x)\ud t\ud x.
  \end{equation}
  The coefficient on $t \ge 0$ depends only on the $H_W$ class of $f_j$.
\end{lemma}

\begin{proof}
  We move derivatives onto the smooth basis test to avoid differentiating the
  singular covariance kernel. We then prove the covariance identity by
  approximation and show that the coefficient depends only on the noise class.

  \medskip
  \noindent\textit{Step 1.~}We prove smoothness of the deterministic
  coefficient.
  Extend $f_j$ smoothly by zero in time and substitute $v = t-u$. The unmasked
  expression is $\int_\R |v|^{-a}\int\Gamma(\ud y)f_j(t-v,x-y)\ud v$. On each
  compact $(t,x)$ set, the integrand involves only bounded $v$ and a common
  compact set of $y$. Local Radon mass and local integrability of $|v|^{-a}$
  provide an integrable majorant for each derivative of the test on these
  fixed sets. The singular kernel is left unchanged under differentiation.
  Thus, every mixed derivative can move onto the smooth test. Dominated
  integration proves smoothness and the derivative bounds. Multiplication by
  $\zeta_R$ gives the stated support. Zero extension is only of a deterministic
  test; no negative-time noise is introduced.

  \medskip
  \noindent\textit{Step 2.~}We identify the coefficient for tests and their
  limits.
  The physical covariance formula and absolute Fubini give
  \eqref{E:driver-literal-covariance} first for real compact tests. The compact
  difference set in $\Gamma$ and integrable temporal singularity justify the
  exchange. For wave tests use the common-support original approximations of
  Proposition~\ref{P:actual-wave-form}, converging in $L^2$ time in the spatial
  supremum norm after masking. The left side converges by noise Cauchy--Schwarz
  and the right side by the coefficient's boundedness on that support. Real
  and
  imaginary parts are treated separately because the original noise is real.
  Combining the two identities gives the complex evaluation defined in the
  preceding section. If two representatives give the same
  class, their Gaussian difference is zero. Testing the covariance identity
  against every compact smooth $F$ makes the difference of their masked
  coefficients zero as a distribution for $t > 0$. Smoothness, then continuity
  at zero, proves the representative assertion on $t \ge 0$. This proves
  Lemma~\ref{L:driver-covariance-coefficients}.
\end{proof}

\subsection{Gain, adjoints and localization}

\begin{lemma}[Deterministic coefficient operators]\label{L:driver-gain-localization}
  For $I \subset [0,T]$, the statewise integral
  \begin{equation}\label{E:driver-coefficient-operator}
    A_j(I)z=\int_I S(-t)B(H_{j,R}(t))S(t)z\ud t
  \end{equation}
  satisfies, for every real $r$,
  \begin{equation}\label{E:driver-gain-one}
    \|A_j(I)\|_{\mathcal{B}(E_r,E_{r+1})}
    \le C_{j,r,R,T}|I|.
  \end{equation}
  Its $E_s$ adjoint has the gain $E_{2s-r-1} \to E_{2s-r}$, so both operators
  act from $E_{s-\gamma}$ to $E$. Choose real compact smooth $\chi = 1$ on
  $B_{2R+T+1}$. On every $E_r$,
  \begin{equation}\label{E:driver-actual-localization}
    D_\chi A_j(I)=A_j(I)=A_j(I)D_\chi.
  \end{equation}
  For each fixed Schwartz pair,
  \begin{equation}\label{E:driver-covariance-row}
    \E[g_j\ell_R(I;z,w)]=\langle A_j(I)z,w\rangle_E.
  \end{equation}
\end{lemma}

\begin{proof}
  We establish the gain in Sobolev order and its adjoint version first. Finite
  propagation then gives both localization identities. Finally, the physical
  pairing identifies the Gaussian coefficient of each row.

  \medskip
  \noindent\textit{Step 1.~}We prove the gain and the adjoint gain.
  The full massless wave is strongly continuous and bounded on every $E_r$ on
  finite time intervals. Its normalized off-diagonal Fourier entries have sizes
  at most $C_T$ and one, respectively; continuity and domination prove the
  strong assertion. Multiplication by $H_{j,R}(t)$ is bounded on every $H^r$,
  continuously in $t$, by weighted Fourier convolution and its compact smooth
  derivative bounds. Hence, $B(H_{j,R}(t)):E_r \to E_{r+1}$ is bounded: its
  second entry is a multiplier of the displacement in $H^r$. The integrand
  in~\eqref{E:driver-coefficient-operator} is strongly continuous with values in
  $E_{r+1}$ for each input in $E_r$, and uniformly bounded in that operator
  norm. Statewise Bochner integration proves~\eqref{E:driver-gain-one},
  including the Lipschitz operator-norm primitive $A_j(0,t)$. Norm continuity of
  $S(t)$ is not needed.

  The zero first component of $B(H_{j,R}(t))z$ belongs to every Sobolev space.
  Its second component lies in $H^r$, exactly the velocity space of $E_{r+1}$.
  The two free propagators preserve their respective Sobolev phase spaces, so
  integration gains one derivative without differentiating the input.
  Under the $E_s$ pairing, the dual of $E_r$ is isometrically $E_{2s-r}$,
  directly by the two Bessel weights. Dualizing the bound with a gain of one
  derivative therefore gives the asserted adjoint gain. It agrees on Schwartz
  phases with the bounded $E_s$ adjoint. Since $\gamma < 1$, gain one at every
  index implies the two stated maps from $E_{s-\gamma}$ to $E$.

  For the adjoint, a functional on the target $E_{r+1}$ is represented under
  the $E_s$ pairing by a vector in $E_{2s-r-1}$. After applying the adjoint it
  is a functional on $E_r$, represented in $E_{2s-r}$. This gives the stated
  indices.

  \medskip
  \noindent\textit{Step 2.~}We prove both localization identities by finite
  propagation.
  For smooth data the free-wave energy on a shrinking ball has nonpositive
  boundary flux: $v\partial_nu-(v^2+|\nabla u|^2)/2 \le 0$. Vanishing initial
  displacement and velocity give zero velocity and gradient in the cone;
  integration of velocity from the zero displacement gives zero displacement
  there as well. Time reversal gives the backward assertion. Thus, both entries
  of $S(\pm t)$ enlarge supports by at most $|t|$. Density and testing against
  compact smooth functions extend this support assertion to Sobolev phases. In
  particular, the forward wave from input vanishing on the larger ball cannot
  reach the support of $H_{j,R}$ by time $T$. For the output, multiplication
  first confines the second entry to $B_{2R}$, and backward propagation
  enlarges that ball by at most $T$. In
  particular, for input $(1-D_\chi)z$, approximate $z$ by Schwartz phases and
  retain the factor $1-D_\chi$; the approximants vanish on $B_{2R+T+1}$. Their
  free displacement vanishes on $B_{2R}$ throughout $[0,T]$, and the same is
  true in the distributional limit. They contribute nothing to the integral. Its
  output is supported in $B_{2R+T}$, by the support of $H_{j,R}$ and backward
  finite propagation. Since $\chi = 1$ on a neighborhood of this ball, both
  identities in~\eqref{E:driver-actual-localization} follow. This uses both
  entries, not a scalar restart or a positivity theorem.

  \medskip
  \noindent\textit{Step 3.~}We identify the covariance row.
  Finally, apply the physical Sobolev pairing
  \eqref{E:variance-physical-adjoint} and~\eqref{E:driver-literal-covariance} to
  the state integral. They give \eqref{E:driver-covariance-row}.
  In~\eqref{E:driver-covariance-coefficient}, the mask is after convolution;
  replacing it by a mask inside $\Gamma*f_j$ would change this identity. This
  proves Lemma~\ref{L:driver-gain-localization}.
\end{proof}

\subsection{Finite variances and first Wiener chaos tails}

Put $J_0(u,v) = (\Lambda^{1/2}u,\Lambda^{-1/2}v)$, $H_0 = L^2(\R^d;\R^2)$, $a_j
= J_0A_j(I)J_0^{-1}$ and $P_\chi = J_0D_\chi J_0^{-1}$. The operator $D_\gamma =
\operatorname{diag}(\Lambda^{-\gamma},\Lambda^{-\gamma})$ acts on these physical
coordinates. Its two compact envelopes $B_{\rm out},B_{\rm in}$ are those of
Proposition~\ref{P:compact-two-envelopes}, after complexification when Schatten
norms are used. For $I \subset [0,T]$ with $0 < h = |I| \le 1$, set
\begin{equation}\label{E:driver-short-variance-budget}
  b_I=C_T h^{2-a}B_\alpha(h),
\end{equation}
where $C_T$ is the uniform-in-$R$ constant in
Proposition~\ref{P:wave-two-variances}.

\begin{lemma}[Both variances and their tails]\label{L:driver-actual-variances}
  For every $N$,
  \begin{equation}\label{E:driver-two-variances}
    V_{{\rm out},N}:=\sum_{j\le N}a_ja_j^*\le b_ID_\gamma^2,
    \qquad V_{{\rm in},N}:=\sum_{j\le N}a_j^*a_j\le b_ID_\gamma^2.
  \end{equation}
  The finite families satisfy both localizations by $P_\chi$. There are bounded
  positive limits $V_{{\rm out},\infty},V_{{\rm in},\infty}$ with positive tails
  $T_{{\rm out},N},T_{{\rm in},N} \to 0$ strongly, and
  \begin{equation}\label{E:driver-compact-tail-envelopes}
    T_{{\rm out},N}\le b_IB_{\rm out},\qquad
    T_{{\rm in},N}\le b_IB_{\rm in}.
  \end{equation}
  For each fixed $z \in E_{s-\gamma}$, the real column maps satisfy
  \begin{equation}\label{E:driver-column-expansions}
    m_{R,I}(z)=\sum_jg_jA_j(I)z,\qquad
    n_{R,I}(z)=\sum_jg_jA_j(I)^{*,E}z
    \quad\text{in }L^2(\Omega;E).
  \end{equation}
\end{lemma}

\begin{proof}
  We first compare finite noise projections with the original covariance
  forms. Summing over a phase basis gives the two variances. We then identify
  the column series and use their positive tails to obtain compact envelopes.

  \medskip
  \noindent\textit{Step 1.~}We project scalar rows onto the noise basis.
  By~\eqref{E:driver-covariance-row}, for real Schwartz $z,w$,
  \[
    \langle X_N(I)z,w\rangle_E
    =\sum_{j\le N}g_j\E[g_j\ell_R(I;z,w)]
    =\mathcal{P}_N\ell_R(I;z,w),
  \]
  where $\mathcal{P}_N$ is the orthogonal projection in $L^2(\Omega)$ onto the
  span of $g_1,\ldots,g_N$. In particular,
  \[
    \sum_{j\le N}|\langle A_j(I)z,w\rangle_E|^2
    =\E|\langle X_N(I)z,w\rangle_E|^2
    \le\E|\ell_R(I;z,w)|^2.
  \]
  The variables $g_j$ are orthonormal in $L^2(\Omega)$. Thus, the squared norm
  of this finite orthogonal projection is the sum of the squared coefficients
  and cannot exceed the squared norm of the original row. This is the
  contraction used in both directions below.
  These are projections onto the noise-test basis, not source regularizations.
  The two variance bounds of Proposition~\ref{P:wave-two-variances}, with the
  mean-square columns of Proposition~\ref{P:wave-square-columns}, apply to these
  same covariance forms and the bound~\eqref{E:driver-short-variance-budget}.

  \medskip
  \noindent\textit{Step 2.~}We sum the row bounds in both phase directions.
  Choose a real Schwartz orthonormal basis $\{e_k\}$ of $E$. Fix real Schwartz
  $w$. For every finite phase truncation $K$,
  \[
    \begin{split}
      \sum_{k\le K}\sum_{j\le N}
      |\langle A_j(I)e_k,w\rangle_E|^2
      &=\sum_{k\le K}\E|\langle X_N(I)e_k,w\rangle_E|^2\\
      &\le\sum_{k\le K}\E|\ell_R(I;e_k,w)|^2
      \le b_I\|w\|_{E_{s-\gamma}}^2,
    \end{split}
  \]
  by~\eqref{E:wave-input-variance}. Let $K \to \infty$. Parseval in the phase
  space and nonnegative monotone convergence, with $N$ still finite, give
  \[
    \sum_{j\le N}\|A_j(I)^{*,E}w\|_E^2
    =\E\|X_N(I)^{*,E}w\|_E^2
    \le b_I\|w\|_{E_{s-\gamma}}^2.
  \]
  Conversely, for fixed real Schwartz $z$, use \eqref{E:wave-output-variance}
  with the output basis varying:
  \[
    \begin{split}
      \sum_{k\le K}\sum_{j\le N}
      |\langle A_j(I)z,e_k\rangle_E|^2
      &=\sum_{k\le K}\E|\langle X_N(I)z,e_k\rangle_E|^2\\
      &\le\sum_{k\le K}\E|\ell_R(I;z,e_k)|^2
      \le b_I\|z\|_{E_{s-\gamma}}^2.
    \end{split}
  \]
  The same phase Parseval limit gives
  \[
    \sum_{j\le N}\|A_j(I)z\|_E^2
    =\E\|X_N(I)z\|_E^2
    \le b_I\|z\|_{E_{s-\gamma}}^2.
  \]
  Thus, both directions use their respective Hilbert-Bessel bounds; neither
  follows from retaining only one variance.

  \medskip
  \noindent\textit{Step 3.~}We extend the finite bounds to the weaker input
  space.
  These inequalities extend at each finite $N$ to the full $E_{s-\gamma}$ input
  space. Indeed, Lemma~\ref{L:driver-gain-localization} gives $A_j:E_{s-\gamma}
  \to E_{s+1-\gamma} \subset E$. Its adjoint gain, at $r = s+\gamma-1$, gives
  the same map for $A_j^{*,E}$. The inclusion holds since $\gamma < 1$.
  Approximate the input by Schwartz phases and pass each of the finitely many
  operator terms to the limit. The resulting constant is still $b_I$,
  independent of $N$; individual $j$-dependent gain constants were not summed.
  Restricting to $E$ and using $\|z\|_{E_{s-\gamma}} = \|D_\gamma J_0z\|_{H_0}$
  gives \eqref{E:driver-two-variances} as two operator quadratic inequalities.
  In physical coordinates, the sum of squared adjoint norms is the quadratic
  form of $\sum_{j\le N}a_ja_j^*$, and the sum of squared direct norms is that
  of $\sum_{j\le N}a_j^*a_j$. The weaker input norm is the quadratic form of
  $D_\gamma^2$. This converts the two scalar inequalities into the displayed
  operator bounds.
  Both localizations follow from~\eqref{E:driver-actual-localization}. No
  self-adjointness or projection property of $P_\chi$ is assumed.

  \medskip
  \noindent\textit{Step 4.~}We identify the first-chaos expansions of the
  columns.
  The coordinates of $m_{R,I}(z)$, for Schwartz $z$, are the original covariance
  forms $\ell_R(I;z,e_k)$ and belong to the closed first Wiener chaos $W(H_W)$.
  Their coefficients against $g_j$ are $\langle A_jz,e_k\rangle_E$. Parseval in
  that chaos, followed by Parseval in $E$ and positive Tonelli, gives
  \[
    \sum_j\|A_jz\|_E^2=\E\|m_{R,I}(z)\|_E^2,
    \qquad m_{R,I}(z)=\sum_jg_jA_jz\quad\text{in }L^2(\Omega;E).
  \]
  For each phase coordinate, completeness of the noise basis gives convergence
  in the scalar first chaos. Summing the squared coefficient tails over the
  phase basis gives the vector-valued mean-square error. Positive Tonelli
  permits this summation, and the finite total column variance forces the
  tails to zero.
  The same argument for the adjoint column gives the other expansion. Membership
  in this first chaos comes from linear noise tests, not merely Gaussianity. To
  extend to a fixed $E_{s-\gamma}$ input, approximate by Schwartz phases. The
  finite inequalities extend to that space by the gain bounds and have a
  constant uniform in $N$. Apply them to the input difference, then use the
  Schwartz expansion and continuity of $m,n$. This
  proves~\eqref{E:driver-column-expansions} on the full claimed input domain.

  \medskip
  \noindent\textit{Step 5.~}We pass to positive variance limits and their
  compact envelopes.
  The convergent positive quadratic sums, restricted to $E$ and transported by
  $J_0$, define the bounded positive operators $V_{\cdot,\infty}$ by
  polarization and Riesz. They retain the bounds $b_ID_\gamma^2$. Their tails
  have quadratic forms tending to zero. Since they are uniformly bounded and
  positive, $\|T_Nv\|^2 \le \|T_N\|\langle T_Nv,v\rangle \to 0$. Passing to the
  limit in the finite localization identities gives
  \[
    V_{{\rm out},\infty}=P_\chi V_{{\rm out},\infty}P_\chi^*,\qquad
    V_{{\rm in},\infty}=P_\chi^*V_{{\rm in},\infty}P_\chi,
  \]
  and the same identities for their tails. Positive congruence applied to $T_N
  \le b_ID_\gamma^2$ proves~\eqref{E:driver-compact-tail-envelopes}.
  The fixed bounded cutoff operators can be passed through these strong
  limits. Subtracting the finite identity from its limiting identity gives the
  same localization for the positive remainder.

  Complexification preserves these statements and constants. The Hilbert norms
  split into real and imaginary squared norms; real operators extend
  isometrically. A real orthonormal basis complexifies to a complex orthonormal
  basis, preserving singular values and multiplicities. The complex column maps
  have the same bounds and sesquilinear row identities. This proves
  Lemma~\ref{L:driver-actual-variances}.
\end{proof}

Lemma~\ref{L:driver-actual-variances} gives tails
$T_{{\rm out},N},T_{{\rm in},N}$ that tend to zero only strongly. They lie
below the compact envelopes in~\eqref{E:driver-compact-tail-envelopes}. We
now use this compact domination to improve the mode of convergence.

\begin{lemma}[Compact domination of positive tails]\label{L:driver-dominated-tails}
  If $0 \le T_N \le B$, with positive compact $B \in S_m$, $m \ge 1$, and $T_N
  \to 0$ weakly, then
  \begin{equation}\label{E:driver-trace-tails}
    \|T_N\|\longrightarrow0,\qquad \operatorname{Tr}T_N^m\longrightarrow0.
  \end{equation}
\end{lemma}

\begin{proof}
  Choose a finite spectral projection $Q$ of $B$ with $\|(1-Q)B(1-Q)\| <
  \epsilon$. Positive quadratic Cauchy--Schwarz gives, for unit $x$,
  \[
    \langle T_Nx,x\rangle
    \le2\langle T_NQx,Qx\rangle
    +2\langle T_N(1-Q)x,(1-Q)x\rangle
    \le2\|QT_NQ\|+2\epsilon.
  \]
  Decompose the unit vector into its $Q$ and $1-Q$ parts. Positive quadratic
  Cauchy--Schwarz bounds the mixed term by twice the geometric mean of the two
  diagonal terms; $2ab \le a^2+b^2$ then gives the factor two in the display.
  The second term is bounded by the spectral tail of $B$.
  Weak convergence makes the finite-dimensional matrix $QT_NQ$ converge in norm
  to zero. Letting $\epsilon\downarrow0$ proves norm convergence.
  Since $T_N$ is positive, its norm is
  the supremum of its quadratic form on unit vectors, so the preceding bound
  proves the norm limit.
  Lemma~\ref{L:compact-positive-order} gives $\lambda_j(T_N) \le \lambda_j(B)$,
  while each $\lambda_j(T_N) \to 0$ by norm convergence. Dominated summation of
  $\lambda_j(B)^m$ proves the trace limit. No order comparison of the operator
  powers is used. This proves Lemma~\ref{L:driver-dominated-tails}.
\end{proof}

Trace powers of variances must still be turned into moments of Gaussian sums.
Proposition~\ref{P:finite-gaussian-matrix} keeps the average of the two trace
powers, but only for matrices. Corollary~\ref{C:finite-gaussian-operators}
bounds compact sums only through the envelope norm $\|K_{\rm out}\|_{S_p}$.

\begin{lemma}[Retained finite Gaussian variance powers]\label{L:driver-gaussian-tail-average}
  Let $Z_\ell$ be finitely many compact operators whose two variances $V_{\rm
  out},V_{\rm in}$ belong to $S_{q/2}$, for an even $q \ge 2$. With independent
  standard real $g_\ell$,
  \begin{equation}\label{E:driver-retained-gaussian-average}
    \E\left\|\sum_\ell g_\ell Z_\ell\right\|_{S_q}^q
    \le(q-1)^{q/2}\frac{\operatorname{Tr}V_{\rm out}^{q/2}
    +\operatorname{Tr}V_{\rm in}^{q/2}}2.
  \end{equation}
\end{lemma}

\begin{proof}
  We retain the two variance trace powers while compressing to finite
  matrices. Compactness then allows passage to the original finite Gaussian
  sum, using Fatou twice.

  \medskip
  \noindent\textit{Step 1.~}We apply the matrix estimate to the compressed
  sum.
  Compress by finite-rank orthogonal $Q_k \to I$ strongly. The compressed output
  variance is at most $Q_kV_{\rm out}Q_k$, since their difference is $\sum_\ell
  Q_kZ_\ell(1-Q_k)Z_\ell^*Q_k \ge 0$; the input inequality is the same with
  adjoints interchanged. Positive-order and compression max-min compare their
  trace powers with the uncompressed ones.
  The comparison of trace powers is made eigenvalue by eigenvalue: positive
  order bounds the compressed variance by the compression of the full
  variance, and max-min bounds the latter eigenvalues by those of the full
  variance. Raising these nonnegative numbers to $q/2$ and summing preserves
  the bounds. Apply
  Proposition~\ref{P:finite-gaussian-matrix} on the finite range, retaining the
  average. For each finite Gaussian coordinate vector, compactness gives
  operator-norm convergence of the compressed sum to the original sum.

  \medskip
  \noindent\textit{Step 2.~}We pass to the uncompressed sum.
  The singular-value norm-Lipschitz argument in the proof of
  Corollary~\ref{C:finite-gaussian-operators}, then Fatou over singular values
  and probability, proves the displayed bound. The extended norm is measurable
  by the same finite-coordinate continuity argument. No infinite Gaussian series
  or $S_q$-Bochner convergence is assumed. This proves
  Lemma~\ref{L:driver-gaussian-tail-average}.
\end{proof}

\subsection{Interval limits on the given probability space}

The two previous lemmas control tails of the series $\sum_jg_jA_j(I)$ through
its two variances. We now fix one interval and use them to pass to the limit in
$N$, without leaving the original probability space.

\LeanMarker{L:driver-interval-limit}
\begin{lemma}[Fixed-interval operator construction]\label{L:driver-interval-limit}
  For each fixed short interval $I$, $X_N(I)$ converges in every finite
  $L^p(\Omega;\mathcal{B}(E))$ to a compact-valued centered Gaussian operator
  $X_R(I)$ on the original probability space. For every admissible even $q \ge
  2$, $q\gamma > d$,
  \begin{equation}\label{E:driver-fixed-schatten}
    (\E\|X_R(I)\|_{S_q}^q)^{1/q}
    \le\sqrt{q-1}\sqrt{C_T}\,h^{1-a/2}B_\alpha(h)^{1/2}
    \|K_\chi\|_{S_q},\qquad K_\chi=P_\chi D_\gamma.
  \end{equation}
  Its fixed Schwartz rows equal the covariance forms in $L^2$.
\end{lemma}

\begin{proof}
  We first show that the finite sums are Cauchy using the compact variance
  tails. We construct their limit on the original probability space, identify
  its rows, and finally extend additively to long intervals.

  \medskip
  \noindent\textit{Step 1.~}We obtain uniform finite moments and vanishing
  tails.
  The finite families $a_j$, $j \le N$, satisfy every hypothesis of
  Proposition~\ref{P:localized-finite-variance} and
  Corollary~\ref{C:finite-gaussian-operators}: both localizations, both weighted
  variances and the independent real $g_j$ have been proved in
  Lemma~\ref{L:driver-actual-variances}. Thus, the coefficients are compact, and
  their finite Gaussian sums satisfy~\eqref{E:driver-fixed-schatten}, uniformly
  in $N$. The fixed cutoff norm may depend on $R,T$.

  Fix an admissible even $q \ge 2$, $q\gamma > d$. Both positive compact
  envelopes in~\eqref{E:driver-compact-tail-envelopes} belong to $S_{q/2}$.
  Lemma~\ref{L:driver-dominated-tails} therefore gives
  \[
    \operatorname{Tr}T_{{\rm out},N}^{q/2}\longrightarrow0,
    \qquad
    \operatorname{Tr}T_{{\rm in},N}^{q/2}\longrightarrow0.
  \]
  For $M > N$, put $Y_{N,M} = \sum_{j=N+1}^Mg_jA_j(I)$. After the unitary $J_0$
  transport, its two finite variances are
  \[
    V_{{\rm out},M}-V_{{\rm out},N}\le T_{{\rm out},N},\qquad
    V_{{\rm in},M}-V_{{\rm in},N}\le T_{{\rm in},N}.
  \]
  Lemma~\ref{L:driver-gaussian-tail-average} and positive eigenvalue max-min
  comparison now give the uniform tail estimate
  \[
    \sup_{M>N}\E\|Y_{N,M}\|_{S_q}^q
    \le(q-1)^{q/2}
    \frac{\operatorname{Tr}T_{{\rm out},N}^{q/2}
    +\operatorname{Tr}T_{{\rm in},N}^{q/2}}2
    \longrightarrow0.
  \]
  For each $M > N$, the two variances of the finite tail lie below the
  corresponding infinite tails. The positive-order comparison bounds their
  trace powers by the right side, which no longer depends on $M$. This is why
  the estimate proves a Cauchy property rather than only boundedness of the
  partial sums.
  The average of the two trace powers is retained. The comparison uses
  eigenvalues, not operator monotonicity of the $q/2$ power. Since $\|Y_{N,M}\|
  \le \|Y_{N,M}\|_{S_q}$, this proves $L^q(\Omega;\mathcal{B}(E))$ Cauchy
  convergence; a uniform moment bound alone would not suffice.

  \medskip
  \noindent\textit{Step 2.~}We construct the operator limit on the prescribed
  probability space.
  All finite sums lie in the separable operator-norm closed real span
  $\mathcal{F}_I$ of the countably many deterministic $A_j(I)$. Choose $N_k$
  increasing so that $\|X_{N_{k+1}}(I)-X_{N_k}(I)\|_{L^q\mathcal{B}(E)} \le
  2^{-k}$. Their $L^1$ norms are summable. Tonelli therefore shows that the
  pointwise operator-norm differences sum absolutely almost surely. Completeness
  gives a limit in $\mathcal{F}_I$ on the original probability space; assign it
  zero on the exceptional null set. It is strongly measurable as a limit of
  separably valued finite Gaussian sums. Fatou and the Cauchy property give
  convergence of the whole sequence in $L^q$ to this limit $X_R(I)$. No new
  noise or probability space has been introduced. Limits constructed at
  different admissible even indices agree by uniqueness of convergence in
  probability. For each requested finite $p \ge 1$, choose an admissible even $q
  \ge p$ and use Lyapunov to obtain $L^p$ operator-norm convergence to that same
  limit. One fixed Schatten index is not used to infer all moment orders.

  Lyapunov compares the $L^p$ norm with the $L^q$ norm on a probability space
  when $p \le q$. Because admissible even $q$ can be chosen arbitrarily large,
  this gives every finite operator-norm moment without changing the random
  limit.

  \medskip
  \noindent\textit{Step 3.~}We identify compactness, moments and Gaussianity
  of the limit.
  The compact operators are norm closed, so an almost surely norm-convergent
  subsequence makes $X_R(I)$ compact valued. Singular-value Fatou applied to the
  finite uniform bound gives \eqref{E:driver-fixed-schatten} at each admissible
  even $q$. Every continuous real linear functional of $X_R(I)$ is an $L^2$
  limit of centered finite Gaussian sums; convergence of means, variances and
  characteristic functions makes it centered Gaussian. The same applies to
  finite linear combinations of these functionals, proving joint Gaussianity on
  the given space.

  \medskip
  \noindent\textit{Step 4.~}We identify the rows as actual noise variables.
  For each fixed real Schwartz $z,w$, the scalar rows $\langle
  X_N(I)z,w\rangle_E$ are the orthogonal projections of the completed covariance
  form $\ell_R(I;z,w)$ onto the complete noise basis. That form lies in the
  original closed first Wiener chaos, so these projections converge to it in
  $L^2(\Omega)$. Operator-norm $L^2$ convergence also gives convergence to
  $\langle X_R(I)z,w\rangle_E$. Uniqueness of the $L^2$ limit identifies these
  two variables on the same space, not merely their laws.

  \medskip
  \noindent\textit{Step 5.~}We extend the construction by finite additivity.
  For a fixed long interval, choose a finite partition into intervals of length
  at most one. Their indicator tests add to the long-interval test except at
  finitely many endpoints. The masked wave products have common compact spatial
  support, and the temporal estimate~\eqref{E:driver-test-continuity} makes
  those endpoint changes have zero original covariance norm. Covariance
  completion and linearity give exact additivity of the forms. The deterministic
  statewise coefficient integrals, hence every $X_N$, add on the same partition.
  Only finitely many short intervals occur in this partition, so the sum of
  their $L^p$ errors tends to zero. No independence of the interval
  contributions is required to pass this finite sum to the limit.
  Passing their finite-sum $L^p$ limits gives the corresponding operator
  additivity and covariance identity for matrix elements, independently of the
  partition by common refinement. These are statements at each fixed interval
  and partition; no uncountable intersection of independently selected null sets
  is used. No long-interval version of the short-$h$ bound is asserted. The
  Schatten conclusion is a fixed-interval moment at admissible even indices, not
  a Schatten-valued path or arbitrary-index claim. This proves
  Lemma~\ref{L:driver-interval-limit}.
\end{proof}

\subsection{Time-path convergence and Gaussian norm bounds}

Lemma~\ref{L:driver-interval-limit} works one interval at a time.
Proposition~\ref{P:actual-operator-driver} asks for more: convergence in the
$C^\theta$ path norm and compact values at every time on one path event.

\LeanMarker{L:driver-holder-limit}
\begin{lemma}[H\"older path limit]\label{L:driver-holder-limit}
  The finite paths $X_N(t) = X_N((0,t))$ converge in
  $L^p(\Omega;C^\theta([0,T];\mathcal{B}(E)))$ for every finite $p \ge 1$, in a
  separable closed path span on the same probability space. The limit starts at
  zero and has compact values for all times on one path event. For sufficiently
  large even $q_0$, its $C^\theta$ norm has $L^q$ bound $C\sqrt q$, uniformly
  over even $q \ge q_0$ and the basis.
\end{lemma}

\begin{proof}
  We first obtain a uniform moment bound on the dyadic increments in a
  stronger Hölder norm. Finite-grid convergence then gives uniform
  convergence, and interpolation gives the desired Hölder convergence. We
  construct the limit in a separable path span.

  \medskip
  \noindent\textit{Step 1.~}We bound increments and the dyadic sum.
  Choose $\theta < \eta < \beta$ and an even $q_0 \ge 2$ with $q_0\gamma > d$
  and $\beta-\eta-1/q_0 > 0$. For $q \ge q_0$, the singular-value sequence gives
  $\|K_\chi\|_{S_q} \le \|K_\chi\|_{S_{q_0}}$: if the latter is nonzero,
  normalize by it, so every singular value is at most one, and sum its $q$th
  powers. The zero case is immediate. Apply the finite Gaussian operator
  estimate with the time-refined variances from
  Proposition~\ref{P:wave-two-variances}. This gives, uniformly in $N$ and even
  $q \ge q_0$,
  \begin{equation}\label{E:driver-increment-moments}
    \|X_N(t)-X_N(u)\|_{L^q(\mathcal{B}(E))}
    \le C\sqrt q\,|t-u|^\beta,\qquad |t-u|\le1.
  \end{equation}
  The constant includes the fixed $S_{q_0}$ cutoff norm and does not depend on
  $q$. Longer increments are finite sums of unit increments, with constants
  depending on $T$.

  Let $D_n(F)$ be the maximum of adjacent increments on the grid $kT2^{-n}$, $0
  \le k \le 2^n$. At fine levels $T2^{-n} \le 1$, the elementary bound for a
  finite maximum yields
  \[
    \|D_n(X_N)\|_{L^q}
    \le 2^{n/q}C\sqrt q\,(T2^{-n})^\beta.
  \]
  The $q$th power of a maximum is at most the sum of the $q$th powers. There
  are $2^n$ adjacent increments, so taking the $q$th root contributes
  $2^{n/q}$. Each increment has length $T2^{-n}$, which yields the displayed
  estimate.
  No independence of these increments is used. The finitely many coarse levels
  are controlled by unit partitions. Minkowski and a geometric sum therefore
  give
  \begin{equation}\label{E:driver-dyadic-moment-sum}
    \sup_N\left\|\sum_{n\ge0}2^{n\eta}D_n(X_N)\right\|_{L^q}
    \le C_\eta\sqrt q,
  \end{equation}
  since the fine-level ratio is at most $2^{-(\beta-\eta-1/q_0)} < 1$.

  \medskip
  \noindent\textit{Step 2.~}We control the stronger Hölder norm by dyadic
  increments.
  For a continuous Banach-valued $F$ starting at zero, this weighted dyadic sum
  controls its $C^\eta$ norm. Approximate two times by dyadic points at a scale
  comparable to their separation. The two grid points differ by at most two
  adjacent increments, and each endpoint's finer binary corrections are bounded
  by the tail sum of $D_n$. Multiplying by the separation to the power $-\eta$
  bounds that total by the weighted sum, up to a constant depending on $T,\eta$.
  Continuity identifies the endpoint limits and also controls the supremum norm
  from $F(0) = 0$. The finite paths are Lipschitz by~\eqref{E:driver-gain-one}.
  Thus, \eqref{E:driver-dyadic-moment-sum} gives uniform $L^qC^\eta$ bounds, and
  the triangle inequality gives the same bound for their differences.

  \medskip
  \noindent\textit{Step 3.~}We prove convergence in the required Hölder norm.
  On a fixed finite time grid, Lemma~\ref{L:driver-interval-limit} gives $L^q$
  convergence of every point value. For continuous $F$,
  \[
    \|F\|_\infty\le\max_k\|F(t_k)\|
    +C(\operatorname{mesh})^\eta\|F\|_{C^\eta}.
  \]
  Apply this to $X_N-X_M$, send $N,M \to \infty$ at a fixed grid, and then send
  the mesh to zero. The uniform $C^\eta$ bound proves sup-norm $L^q$ Cauchy
  convergence. The deterministic interpolation inequality
  \begin{equation}\label{E:driver-holder-interpolation}
    \|F\|_{C^\theta}
    \le C\|F\|_\infty^{1-\theta/\eta}
    \|F\|_{C^\eta}^{\theta/\eta}
  \end{equation}
  follows by bounding each increment by both $2\|F\|_\infty$ and its $C^\eta$
  bound, and splitting at their balancing scale. Its supremum term is controlled
  too. Scalar H\"older in probability proves the $L^qC^\theta$ Cauchy property.
  In the interpolation estimate, apply scalar Hölder to the two random factors
  with reciprocal exponents $1-\theta/\eta$ and $\theta/\eta$. This bounds the
  $L^qC^\theta$ norm by the corresponding powers of the $L^q$ supremum norm
  and the $L^qC^\eta$ norm. The first tends to zero for differences and the
  second stays bounded.
  For any requested finite $p$, choose an even $q \ge \max(p,q_0)$ and use
  Lyapunov.

  \medskip
  \noindent\textit{Step 4.~}We choose the measurable limit and one path event.
  The full H\"older/operator space need not be separable. Nevertheless, every
  finite path lies in the span of the countably many deterministic Lipschitz
  paths $A_j(0,\cdot)$. Their $C^\theta$-closed real span is a separable Banach
  space. In it the same elementary Cauchy/subsequence construction as in
  Lemma~\ref{L:driver-interval-limit} gives a strongly measurable limit $M_R$ on
  the original space. Different $p$ limits agree by uniqueness of convergence in
  probability in the path norm. An almost surely $C^\theta$-convergent
  subsequence converges uniformly in operator norm. Finite paths have compact
  values on short intervals, and on long ones by finite additivity. Therefore,
  the limit is a norm limit of compact operators at every time on this same
  event. Uniform convergence is essential: one does not select a separate
  exceptional set for each time. Thus,
  the limit has compact values at every time on that single path set. It starts
  at zero. Passing the uniform bounds from~\eqref{E:driver-dyadic-moment-sum}
  through the path convergence proves the asserted $C\sqrt q$ bound. This proves
  Lemma~\ref{L:driver-holder-limit}.
\end{proof}

\begin{proof}[Proof of Proposition~\ref{P:actual-operator-driver}]
  We assemble the interval and path constructions, then verify basis
  independence. The uniform square-root moment bound gives the exponential
  norm estimate by summing its power series.

  \medskip
  \noindent\textit{Step 1.~}We identify the path and its covariance rows.
  Lemmas~\ref{L:driver-noise-basis}--\ref{L:driver-holder-limit} construct the
  stated finite sums and their path limit, with no new probability space.
  Evaluation at a fixed time agrees with the interval limit by uniqueness in
  probability. Its covariance identity for matrix elements is
  \eqref{E:driver-raw-row}, by Lemma~\ref{L:driver-interval-limit}, including
  the long-interval interpretation. The same additivity identifies each fixed
  increment with $X_R(I)$, proving its stated Schatten moment bound.

  The finite paths are Gaussian linear combinations of deterministic paths.
  Every continuous linear functional on their closed path span converges in
  $L^2$ to a centered Gaussian; finite linear combinations do as well. Thus,
  $M_R$ is a centered Gaussian path. The four real pair identities give its
  complexified sesquilinear row identity.

  \medskip
  \noindent\textit{Step 2.~}We compare paths built from different noise bases.
  The construction used only compact smooth representatives, orthonormality and
  completeness, so it applies to any such noise basis. Compare two constructions
  using a fixed countable dense Schwartz phase collection and rational times.
  Both rows equal the same covariance forms. A countable null-set union,
  boundedness of the operators, phase density and path continuity give equality
  at all times on one probability-one set. This proves indistinguishability
  without an assumed common event for uncountably many original evaluations.

  \medskip
  \noindent\textit{Step 3.~}We deduce all finite moments and exponential
  integrability.
  For $p \ge 1$, choose even $q \ge \max(p,q_0)$ with $q \le p+q_0+2 \le
  (q_0+3)p$. Lemma~\ref{L:driver-holder-limit} and Lyapunov give the first bound
  in~\eqref{E:driver-gaussian-norm}. Write its resulting constant as $C_1$;
  Lyapunov and $q \le (q_0+3)p$
  absorb the passage from even $q$ into the constant $C_1$. Now
  let $Y = \|M_R\|_{C^\theta}$. If $C_1 = 0$, then $Y = 0$ almost surely and the
  exponential bound is immediate. Otherwise, $\E Y^{2m} \le
  (C_1\sqrt{2m})^{2m}$. Since $m! \ge (m/e)^m$, as follows by integrating $\log
  x$ from one to $m$, the $m$th term in the exponential series is at most
  $(2ecC_1^2)^m$. For $0 < c < 1/(2eC_1^2)$, positive Tonelli and the geometric
  series prove the second bound. This proves
  Proposition~\ref{P:actual-operator-driver}.
\end{proof}

All constants here may depend on $R,T$ and the fixed admissible parameters. No
uniform-in-$R$ path-norm estimate or $C^\theta S_q$-valued path is asserted.
This constructs the local masked integrated operator on the original noise.
Source-cutoff convergence, Young evolution, canonical initial data and
solution-moment asymptotics remain separate parts of the argument.
